% !TeX root = ../main.tex
\section{Universal simulation}
\label{sec:universal-simulation}

To simulate circuits that leave the DCP state class, we represent the evolving state as a linear combination of DCP states. The analogous construction using Pauli stabilizer states is quantified by stabilizer rank.

\begin{definition}[Stabiliser rank \cite{bravyi2016improved}]\label{def:stabilizer-rank}
    The \emph{stabilizer rank} $\chi(\psi)$ of a state $\ket{\psi}$ is the smallest integer $\chi$ such that $\ket{\psi}$ can be written as a linear combination $\ket{\psi} = \sum_{i=1}^\chi c_i\ket{\varphi_i}$ of stabiliser states $\ket{\varphi_i}$ with complex coefficients $c_i$.
\end{definition}

The notion of stabiliser rank easily generalises to \emph{DCP rank}:

\begin{definition}[DCP rank]\label{def:dcp-rank}
    The \emph{DCP rank} $\zeta(\psi)$ of a state $\ket{\psi}$ is the smallest integer $\zeta$ such that $\ket{\psi}$ can be written as a linear combination $\ket{\psi} = \sum_{i=1}^\zeta c_i\ket{\varphi_i}$ of DCP states $\ket{\varphi_i}$ with complex coefficients $c_i$.
\end{definition}

Since stabiliser states are DCP states, we trivially have $\zeta(\psi) \leq \chi(\psi)$.
Furthermore, if $\ket{\psi}$ is prepared by a $\CNOT+\T+\Hgate$ circuit with $h$ Hadamards, then we trivially have $\zeta(\psi) \leq 2^h$.
This is because each Hadamard gate can be decomposed into $\Hgate = (X + Z) / \sqrt 2$, leaving only gates that preserve DCP states.
In general, we expect $\zeta$ to scale exponentially, however there might be decomposition strategies that achieve better scaling than this naive approach.
We leave it to future work to investigate this question and improve upon the bounds given here.

Given an explicit decomposition
\begin{align}
    \ket{\psi}=\sum_{j=1}^{R}c_j\ket{\phi_j}
\end{align}
into DCP states, we can compute any computational-basis amplitude by evaluating the amplitude of each term and summing the contributions.
By \cref{thm:efficient-simulation}, this requires $O(n^3R)$ arithmetic operations.
Here $R$ denotes the number of terms in the supplied decomposition, which need not achieve the minimum DCP rank $\zeta(\psi)$.
The cost of constructing the decomposition must also be accounted for.
Throughout, the global-phase convention of each DCP term must be tracked consistently with its coefficient $c_j$.
Weak simulation, that is, sampling computational-basis outcomes, requires an additional argument.
Computing marginal probabilities directly introduces $R^2$ cross terms involving overlaps of projected DCP states, which our polynomial-time algorithms do not generally evaluate.
Several approaches have been proposed in the stabilizer-decomposition literature to avoid or mitigate analogous difficulties.
For example, Refs.~\cite{bravyi2016improved, bravyi2019simulation} approximate marginal probabilities via a norm estimation technique based on overlaps with carefully chosen random stabiliser states.
Ref.~\cite{bravyi2019simulation} also proposed a Metropolis-Hastings Markov Chain Monte Carlo method for approximate sampling.
Finally, Bravyi and Liu \cite{bravyi2022how} introduced the \emph{gate-by-gate} algorithm for exact sampling without computing marginals.
The gate-by-gate algorithm applies directly because it requires only computational-basis amplitudes of circuit prefixes.
We describe this construction below.

Suppose $\ket{\psi} = C\ket{0}$ is prepared by some $\CNOT+\T+\Hgate$ circuit $C = U_m \cdots U_1$ made up of gates $U_i$.
Our goal is to sample from the distribution $P(x) := |\bra{x}C\ket{0}|^2$.
Let $C_k = U_k \cdots U_1$ be the subcircuit containing only the first $k$ gates and define the output distribution at time step $k$ as $P_k(x) := |\bra{x}C_k\ket{0}|^2$.
The idea behind gate-by-gate sampling is to start with a sample from the trivial distribution $P_0$ (i.e. the all-zero bitstring) and update it one gate at a time until we arrive at a sample that is distributed according to $P_m = P$.
This is because adjacent distributions $P_{k-1}$ and $P_k$ are closely related since they only differ by a single gate $U_k$.
Suppose $U_k$ is a $\T$ gate on qubit $i$, then
\[ P_k(x) = |\bra{x}\T_iC_{k-1}\ket{0}|^2 = |\bra{x}C_{k-1}\ket{0}|^2 = P_{k-1}(x), \]
i.e. the distribution stays the same.
For a $\CNOT$ between qubits $i$ and $j$ we have
\begin{align}
    \begin{split}
        P_k(x) &= |\bra{x}\CNOT_{ij}C_{k-1}\ket{0}|^2 \\
        &= |\bra{x \oplus x_ie_j}C_{k-1}\ket{0}|^2 \\
        &= P_{k-1}(x \oplus x_ie_j),
    \end{split}
\end{align}
so a sample $x \sim P_{k-1}$ can be updated to $x' := x \oplus x_ie_j \sim P_k$.
Applying a Hadamard on the other hand will require some more work.
Without loss of generality, suppose $U_k$ is a Hadamard on the first qubit.
Writing $y=(x_2,\ldots,x_n)$, we have
\begin{align}
    P_k(x_1,y)
    &=P_k(x_1\mid y)P_k(y)\\
    &=P_k(x_1\mid y)P_{k-1}(y)\\
    &=\frac{P_k(x_1,y)}{P_k(0,y)+P_k(1,y)}P_{k-1}(y),
\end{align}
whenever $P_{k-1}(y)>0$.
Here $P_k(y)=P_{k-1}(y)$ because the Hadamard acts only on the first qubit.

Thus, given a sample $(x_1,y)\sim P_{k-1}$, we retain $y$ and draw a new first bit according to
\begin{align}
    \Pr(x'_1=b\mid y)
    =\frac{P_k(b,y)}{P_k(0,y)+P_k(1,y)},
    \qquad b\in\{0,1\}.
\end{align}
The denominator is positive for every $y$ sampled with nonzero probability.
Each Hadamard therefore requires two computational-basis amplitude evaluations, whose squared magnitudes give the required probabilities.
A circuit containing $h$ Hadamards consequently requires $2h$ such evaluations.
Let $R_j$ be the number of terms in a DCP decomposition of the state immediately after the $j$th Hadamard.
Excluding the cost of constructing these decompositions and sampling from the resulting probabilities, the amplitude evaluations required for one sample take
\begin{align}
    O\left(n^3\sum_{j=1}^{h}R_j\right)
\end{align}
arithmetic operations.
When using the previously discussed decomposition $H=(X+Z)/\sqrt{2}$, we get $R_j := 2^j$ such that
\begin{align}
    \sum_{j=1}^{h}R_j
    =\sum_{j=1}^{h}2^j
    =2^{h+1}-2
    = O(R_h).
\end{align}
In other words, the cost of all intermediate amplitude evaluations is asymptotically dominated by the cost of evaluating the final state with $R_h$ terms.
This means that gate-by-gate sampling matches the complexity of strong simulation in that setting.
The same conclusion holds for any decomposition strategy whose prefix term counts satisfy $\sum_{j=1}^{h}R_j=O(R_h)$.
